%% 段組ごとの段間と段間罫（\evenendcolumns[columnsep=...,columnseprule=...]{段数}）
\documentclass[twocolumn,paper=b5,fontsize=9pt]{jlreq}
\usepackage[debug]{evenend}
\setlength\columnseprule{0.4pt}
\newcommand\para{吾輩は猫である．名前はまだ無い．どこで生れたかとんと見当がつかぬ．何でも薄暗いじめじめした所でニャーニャー泣いていた事だけは記憶している．\par}
\makeatletter
\newcommand\showsep{\typeout{BLOCK N=\the\evenend@N\space sep=\the\evenend@colsep\space rule=\the\evenend@colrule\space width=\the\columnwidth}}
\makeatother
\begin{document}
\section{既定}\showsep
\para\para
\evenendcolumns[columnsep=1\zw,columnseprule=0pt]{3}
\section{段間1字，罫なし}\showsep
\para\para\para
\evenendcolumns[columnsep=4\zw,columnseprule=2pt]{2}
\section{段間4字，太い罫}\showsep
\para\para
\evenendcolumns{2}
\section{既定に戻る}\showsep
\para\para
\setlength\columnsep{3\zw}
\begin{multicols}{3}
\showsep
\para\para\para
\end{multicols}
\showsep
\setlength\columnsep{2\zw}
\evenendcolumns[columnseprule=1pt]{2}
\section{ページをまたぐ}\showsep
\para\para\para\para\para\para\para\para\para\para\para\para\para\para\para\para\para\para\para\para
\end{document}
